\begin{textsample}{POS Dim 3 – ce – Score 60.00 – ce/ce\_inf\_7.txt}  \label{ex:f3_pos_051}
4.1.1.5.
Aprender a partir das interações e da \textbf{brincadeira} O processo educativo, que visa a possibilitar condições adequadas de educação e \textbf{cuidados}, requer a garantia do respeito ao ritmo próprio de cada \textbf{criança}, além da compreensão de que ela aprende desde o \textbf{nascimento}, nas experiências cotidianas vividas, por meio do corpo, da emoção, da linguagem verbal, das interações e das múltiplas linguagens.
Aprender pode ser entendido como o processo de modificação do modo de agir, sentir e pensar de cada pessoa, e que não pode ser atribuído à mera maturação orgânica, mas à construção de sentidos e significados sobre si mesmo, sobre os outros e sobre o mundo que se dá pela multiplicidade das experiências vividas.
Nessa concepção, as possibilidades de aprendizagem não são resultado de processos espontâneos, mas requerem alguns elementos mediadores, em especial, a colaboração de diferentes \textbf{parceiros} na realização de alguma tarefa.
O ensinar pode ser compreendido como um processo dinâmico, interativo, de inúmeras possibilidades que contribuem, por meio de experiências inaugurais ou ressignificadas, à inserção dos sujeitos ao meio sociocultural.
Esse processo se concretiza em situações diversas e cotidianas.
Por exemplo, quando se diz para um \textbf{bebê} que é hora de mamar e que ele deve estar com fome, ou quando é chamado pelo seu nome, ou ainda quando \textbf{escuta} o nome das pessoas que circulam em seu meio.
Esses significados culturais poderão ser apropriados a partir das experiências desses sujeitos, por isso se ressalta que é na interação com o outro que se dá significado às coisas.
Fica mais fácil compreender essa ideia quando se pensa em situações que são vividas pelas pessoas e quando cada uma relata o fato interpretado, a partir do seu ponto de vista.
Essa perspectiva para a Educação \textbf{Infantil} requer aportes teóricos e práticos, consistentemente fundamentados nas necessidades e características de \textbf{crescimento} e desenvolvimento de \textbf{bebês}, \textbf{crianças} bem \textbf{pequenas} e \textbf{crianças} \textbf{pequenas}.
A presença de professoras, de professores que se portam de forma \textbf{atenta}, sensível, respeitosa, mediadora das valiosas e significativas aprendizagens das \textbf{crianças}, que instigam seus \textbf{afetos}, suas sensações, seus movimentos, possibilitando as interações e as \textbf{brincadeiras}, sua memória, linguagem e identidade, ampliarão as possibilidades de construção de conhecimentos já sistematizados na cultura, e, especialmente, dos sentidos pessoais que decorrem das interpretações \textbf{infantis}.
Mesmo quando uma \textbf{professora}, um professor propõe que as \textbf{crianças} se envolvam em uma atividade em princípio similar, como por exemplo, o desenho ; se houver liberdade de ação e interação pelas \textbf{crianças} na criação desse desenho — o que é demandado quando se lê com atenção as DCNEI e a BNCC — cada \textbf{criança} produzirá graficamente o suporte com singularidade, lançando \textbf{mão} de suas experiências e conhecimentos prévios, interagindo com aquela situação, constituída naquele momento.
Algumas \textbf{crianças} fazem o desenho em silêncio, outras precisam \textbf{conversar} com o \textbf{colega} do lado, outras se levantam da sua cadeira e ficam em \textbf{pé} para experimentar diferentes possibilidades de produção.
Enfim, a proposta de desenhar pode até ser feita coletivamente, mas as experiências das \textbf{crianças} são as mais diversas.
Essa diversidade só enriquece o ambiente de aprendizagem.
Dessa forma, ações de ensino e aprendizagem podem partir de \textbf{adultos}, \textbf{crianças} e situações interativas do cotidiano, mas não se esgotam nesses elementos.
Em outras palavras, não se aprende só com a \textbf{professora}, o professor, mas com diferentes elementos simbólicos que agem como recursos na relação da \textbf{criança} com o mundo. - amplia o olhar para as diferentes fontes de aprendizagem ( \textbf{adultos}, \textbf{crianças} e situações diversificadas de interação ) ; - opõe -se à ideia de ensino como movimento que parte apenas da \textbf{professora}, do professor e que toma a \textbf{criança} como mero receptor de suas mensagens ; - considera que o processo de ensino e aprendizagem depende da interação que se estabelece entre \textbf{adultos} e \textbf{crianças}, e suas decorrências dependerão basicamente da atividade de cada \textbf{criança}, que continuamente atribui sentidos aos significados que lhe são apresentados, sem que esse reconhecimento enfraqueça a importância das ações da \textbf{professora}, do professor.
Com base nessa concepção, propõe -se pensar qual é o papel do professor.
Para tanto, tem -se que afastar a ideia da \textbf{professora}, do professor, tradicionalmente, associado a alguém que é um transmissor de conhecimentos às \textbf{crianças}.
O desafio é buscar uma nova forma de pensar como o professor deve atuar junto às/com as \textbf{crianças} \textbf{pequenas}, bem \textbf{pequenas} e junto aos/com os \textbf{bebês} em articulação com as famílias e os demais profissionais da escola de educação \textbf{infantil}.
No caso da \textbf{instituição} educacional, o professor é um mediador das aprendizagens das \textbf{crianças}.
Ele atua de modo : - indireto, pelo arranjo do ambiente de aprendizagem das \textbf{crianças}, onde outros mediadores estão presentes : os espaços, os objetos, as indumentárias, os \textbf{livros}, os horários, os agrupamentos \textbf{infantis} ; - direto, quando, por exemplo, ao interagir com as \textbf{crianças}, lhes apresenta novidades, dando suporte para elaboração, validação ou negação das hipóteses ; quando as pega no \textbf{colo} numa situação de insegurança ; quando medeia normas de conduta, e outras ações.
Além das interações dos \textbf{adultos} com as \textbf{crianças}, a BNCC ( BRASIL, 2017 ) ressalta também as interações das \textbf{crianças} com outras \textbf{crianças}, que são \textbf{parceiras} na fascinante tarefa de compreender o mundo e a si mesmas.
As \textbf{crianças}, nas interações que estabelecem entre si, aprendem a : - fazer amigos, negociar significados e decisões, resolver conflitos, partilhar sentimentos e combater estereótipos e preconceitos que limitam o desenvolvimento de uma pessoa ; - viver em grupo, a ser sensível ao ponto de vista ou aos sentimentos do outro, a cooperar em diferentes tarefas, a conhecer suas limitações e possibilidades, a aceitar -se e aos companheiros, a controlar seus impulsos e emoções, e a desenvolver variadas formas de comunicação e de expressão de \textbf{afetos} ; - partilhar com outras \textbf{crianças} os conhecimentos e a identidade que o grupo lhes oferece ; - construir sua imagem pessoal no grupo.
Tanto as interações da/do professora/professor com as \textbf{crianças}, das \textbf{crianças} entre si quanto às interações com os espaços e materiais estimulam processos de aprendizagem que fazem avançar o desenvolvimento.
Em uma situação, as explicações dadas pela/pelo professora/professor, ou a narrativa de um caso por outra \textbf{criança}, possibilitam a aprendizagem.
Em outra ocasião, a \textbf{escuta} de uma história, o folheio de um \textbf{livro}, a participação em um faz de conta, a construção de um castelo com sucata, e outras atividades, são poderosos mediadores da conquista pela \textbf{criança} de novas formas de agir, pensar e sentir.
As interações das \textbf{crianças} As interações sociais constituem meios para os processos de aprendizagem e desenvolvimento humano, já que, em seu acontecimento, colocam em jogo produções e ferramentas da cultura.
É por essa importância que elas são consideradas como um dos eixos do trabalho pedagógico dedicado aos \textbf{bebês} e \textbf{crianças}, nos espaços educacionais coletivos.
Na Educação \textbf{Infantil}, sob a ótica das \textbf{crianças}, essas interações podem ocorrer entre : - As \textbf{crianças} e as/os professoras/professores adultas/adultos — essenciais para dar riqueza e complexidade às \textbf{brincadeiras} e demais ações promovidas no cotidiano, seja na \textbf{instituição} de educação \textbf{infantil} e em diferentes contextos de socialização da \textbf{criança} ; - \textbf{Crianças} e \textbf{crianças} — permite o compartilhamento de experiências e informações entre as \textbf{crianças}, propiciando ampliarem ou construírem conhecimentos e se desenvolverem na relação eu — outro.
Portanto, as interações sociais ocorrem nas relações das \textbf{crianças} com o \textbf{adulto}, da \textbf{criança} com a \textbf{criança}, da \textbf{criança} com a cultura, por meio dos processos de mediação simbólica, existentes nos objetos, nas expressões, \textbf{gestos}, na \textbf{brincadeira}, nos ambientes, na palavra, e nas manifestações das múltiplas linguagens da \textbf{criança} ( ARAÚJO, 2017 ).
Especificamente nas \textbf{instituições} de Educação \textbf{Infantil}, é esperado que as interações ocorram em situações concretas e cotidianas, significativas, transformadoras dos espaços em ambientes vivos e dinâmicos, nos quais as \textbf{crianças} se revelam como protagonistas, questionam, dialogam, convivem e participam, construindo sua identidade individual e coletiva.
Tal possibilidade diverge de uma perspectiva na qual a \textbf{criança} só aprende ao \textbf{escutar} atentamente o \textbf{adulto}.
Ao proporcionar diversas experiências, sejam livres ou dirigidas e desafiantes, tais ações contribuem para que as \textbf{crianças} expressem de diversas maneiras seus sentimentos, argumentos, hipóteses, saberes e atuem ativamente na cultura e, com os seus \textbf{parceiros}, produzam as culturas \textbf{infantis}.
Vale ressaltar ainda que existem situações em que “ [... ] as interações professora/professor-criança são mais necessárias ” uma vez que este é profissional de referência para a \textbf{criança}, \textbf{apoiando} as iniciativas \textbf{infantis}, \textbf{acolhendo} os medos, as conquistas, as \textbf{curiosidades}, as alegrias e questionamentos.
A interação qualitativa implica não apenas promover boas experiências para as \textbf{crianças}, mas significa acolhê -las desde que ingressem na \textbf{instituição}, promover o seu bem-estar, considerar o que elas dizem, “ [... ] perceber e respeitar suas características, compreender seus motivos e estabelecer vínculos afetivos [... ] ” ( SEDUC, 2011, p. 38 ).
Enquanto em outras situações, as interações das \textbf{crianças} são os recursos mais importantes para as aprendizagens, especialmente durante as \textbf{brincadeiras}, ocasiões em que realizam imitações, opõem -se, disputam objetos e \textbf{brinquedos}, combinam ações, criam enredos e colocam em jogo novas ideias, improvisam falas, entre outras.
Para que as \textbf{partilhas} entre as \textbf{crianças} ocorram de forma democrática, com segurança afetiva, é necessário que o professor planeje as ações, tendo a participação das \textbf{crianças} nas decisões sobre os espaços, sobre os materiais pedagógicos, \textbf{livros}, \textbf{brinquedos} etc., por meio de experiências que possibilitem a \textbf{escuta} e expressão das múltiplas linguagens.
Por meio de práticas interativas e significativas que fomentem a democracia e a participação, garantindo às \textbf{crianças} vivenciar os princípios éticos, políticos e estéticos, a/o professora/professor de Educação \textbf{Infantil} promoverá o respeito aos direitos de cada \textbf{criança}, o que difere “ [... ] das formas de trabalho que buscam disciplinar as \textbf{crianças} por castigos, ameaças e outras formas de violação dos direitos \textbf{infantis} à proteção, ao \textbf{cuidado} e à educação. ” ( Ibidem ).
Na BNCC, as interações das \textbf{crianças} com as outras \textbf{crianças} e com os \textbf{adultos} são eixo fundamental da organização do trabalho pedagógico.
Os Direitos de Aprendizagem e os Campos de Experiências — e em cada Campo, os Objetivos de Aprendizagem e Desenvolvimento — estão sustentados na ideia de que a interação promove novas e diversificadas aprendizagens.
A \textbf{brincadeira} como atividade da \textbf{criança} É por meio da \textbf{brincadeira} que as \textbf{crianças} constituem a si mesmas, dão significado ao mundo ao seu redor, expressam seus \textbf{desejos}, medos, experimentam papéis, interagem com seus \textbf{pares} e compartilham suas culturas.
O \textbf{brincar} é inerente à cultura das \textbf{crianças} e a Educação \textbf{Infantil} possui um importante papel ao garantir esse direito e ampliar esse repertório que elas já possuem, introduzindo -as nos mais variados tipos de \textbf{brincadeiras}.
Parecer da CNE/CP nº 15/2017, que cita a Resolução ( CNE/CEB nº 5/2009 ), assegura o \textbf{brincar} como um dos direitos de aprendizagem e desenvolvimento da \textbf{criança}, no âmbito da Educação \textbf{Infantil}, ressaltando que é direito da \textbf{criança} : \textbf{brincar} cotidianamente de diversas formas, em diferentes espaços e tempos, com diferentes \textbf{parceiros} ( \textbf{crianças} e \textbf{adultos} ), ampliando e diversificando seu acesso a produções culturais, seus conhecimentos, sua imaginação, sua criatividade, suas experiências emocionais, corporais, \textbf{sensoriais}, expressivas, cognitivas, sociais e \textbf{relacionais} ( INCISO II, p. 53 ).
É importante destacar que, ao garantir o \textbf{BRINCAR} em sua perspectiva mais ampla e diversificada, o projeto pedagógico também caminha para a promoção dos demais Direitos de Aprendizagem e Desenvolvimento.
Afinal, ao \textbf{BRINCAR} as \textbf{crianças} aprendem a CONVIVER ; a PARTICIPAR ativamente em seus grupos de \textbf{parceiros} ; a EXPLORAR linguagens, materiais, espaços, elementos da cultura ; a EXPRESSAR emoções, sentimentos, criações, ideias, hipóteses etc. e, por fim, a CONHECER -SE, já que na interação com o outro, na elaboração das \textbf{brincadeiras} e realização dos jogos, compreende as suas necessidades, ideias, \textbf{afetos}, ao compartilhar \textbf{afetos} e ideias diversas que a colocam em outra perspectiva.
Por essa abrangência, a \textbf{brincadeira} é considerada como um dos eixos do trabalho pedagógico na Educação \textbf{Infantil}.
Promover situações na Educação \textbf{Infantil} por meio de práticas pedagógicas que permitam que as \textbf{crianças} \textbf{brinquem} sozinhas, com seus \textbf{pares}, \textbf{brincadeiras} livres ou dirigidas, é uma forma de garantir sua cidadania.
É importante salientar que a \textbf{criança} não nasce sabendo \textbf{brincar}, assim, essa ação é aprendida por meio das interações crianças/adultos, criança/criança, ou no contato com diferentes materiais e artefatos ( \textbf{brinquedos}, \textbf{livros}, materiais não estruturados, dentre outros ).
Na observação de outras \textbf{crianças} ou quando a \textbf{professora} propõe \textbf{brincadeiras} compostas por regras, as \textbf{crianças} vão aprendendo a criar e a reproduzir a ação do \textbf{brincar}.
Ao participar de \textbf{brincadeiras} com \textbf{adultos} e com outras \textbf{crianças} mais experientes, ela vai aprendendo e se apropriando das regras, enquanto \textbf{brinca}.
Basta ver como as \textbf{crianças} \textbf{brincam} no início do ano e o quanto suas \textbf{brincadeiras} vão ganhando novos elementos no decorrer da convivência e das diversas experiências que vive na Educação \textbf{Infantil}.
Para isso, o trabalho pedagógico deve ser rico em linguagens e manifestações artístico-culturais e em situações de observação do mundo e das relações sociais.
No jogo simbólico da \textbf{brincadeira}, as \textbf{crianças} apropriam -se de informações da cultura do mundo \textbf{adulto}, reportando -se a situações vivenciadas por elas em seus contextos sociais, no entanto não se limitam a interiorizá -la, mas, para além disso, interpretam -na criativamente num processo denominado pelo sociólogo William Corsaro ( 2011 ) de “ reprodução interpretativa ”, uma vez que ao reproduzir uma ação vivenciada, a \textbf{criança} não faz uma \textbf{cópia} daquilo que vivenciou, mas a interpreta a partir da sua leitura do vivido na busca por lidar com inquietações próprias e exclusivas, experimentadas nos modos de ser e de se relacionar com as outras \textbf{crianças}.
Nesse sentido, é \textbf{interessante} enxergar em uma \textbf{brincadeira} de papéis registrada em uma escola de educação \textbf{infantil}, localizada em uma comunidade ribeirinha da Amazônia, os elementos da cultura, do modo de viver da comunidade, compondo a \textbf{brincadeira} das \textbf{crianças}.
Episódio : TOMANDO BANHO NA MARÉ⁶ Participantes : Clara ( 5,013 ), Lucas ( 5 ; 2 ), Marcos ( 4 ; 1 ) e Paulo ( 5,0 ) Lucas, Clara e Marcos \textbf{brincam} em um canto da sala.
Clara está com uma boneca na \textbf{mão} e os \textbf{meninos} com alguns \textbf{brinquedos}, como barcos e peças de um jogo de \textbf{montar}.
Clara inicia um diálogo com os \textbf{meninos}.
Clara — Meu filho já sabe subir na árvore.
Ele sobe e não cai.
Lucas — Não cai?
Então ele já sabe.
Mas ele ainda não sabe tomar banho na maré.
Clara — Sabe.
Lucas — Então, vamos lá no igarapé ver se ele sabe.
Clara — Vamos.
Clara — ( Levanta -se com a boneca na \textbf{mão} e acompanha Lucas, que se dirige a um outro canto da sala.
Marcos e Paulo também o acompanham ).
Marcos — Ele vai é morrer afogado.
Clara — Não vai não.
Ele já sabe nadar.
Marcos — Ela diz que o filho dela sabe tudo. ( As \textbf{crianças} chegam no local proposto por Lucas e ele fala ).
Lucas — Isso que eu quero ver.
Pronto, põe ele na água.
Vamos ver se ele sabe tomar banho na maré. ( Clara pega a boneca e vai colocando bem devagar no chão, como se a estivesse colocando na água ).
Lucas — Agora vem a água da maré.
Olha a água chegando.
Está ficando forte, forte. ( Todos olham para a boneca imóvel no chão ).
Lucas — É.
O filho dela parece que não tem medo.
Acho que ele já sabe tomar banho na maré.
Em um breve episódio de \textbf{brincadeira} simbólica, é possível ver todo o empreendimento das \textbf{crianças} na retomada de uma situação do cotidiano e recriação no espaço da escola.
Pensamento, fala, memória são movimentados pelas \textbf{crianças} que precisam reconstruir uma cena de beira de rio e maré, que está presente em imagem mental que é colocada em ação. \textbf{Gestos} entram em cena com toda a expressividade das \textbf{crianças}, que assumem papéis e precisam ser convincentes umas com as outras e \textbf{consigo} mesmas.
Atividades da \textbf{criança} que, na \textbf{brincadeira}, a impulsionam a ir além do seu desenvolvimento, do seu cotidiano.
Percebem como as mães da comunidade lidam com seus filhos \textbf{pequenos} e valorizam determinados aprendizados.
Os outros \textbf{colegas} a desafiam como “ bons vizinhos ” que são!
E assim, a \textbf{brincadeira} segue reproduzindo significados da cultura e criando novos sentidos para a \textbf{brincadeira} das \textbf{crianças}.
Situações como essa acontecem cotidianamente na Educação \textbf{Infantil}.
Basta parar e registrar momentos de \textbf{brincadeira} das \textbf{crianças} nos diferentes tempos de sua jornada.
Será fácil perceber o que mais chama a atenção delas nas práticas da cultura que constituem a amplitude do estado do Ceará e a diversidade dos municípios cearenses.
A liberdade de escolha e o protagonismo da \textbf{criança} são essenciais na \textbf{brincadeira}.
O caráter de liberdade se expressa pelo fato de a \textbf{criança} poder escolher, mediante a sua \textbf{vontade}, o tema da \textbf{brincadeira} com a qual quer se envolver e o papel que desempenhará dentro dela, entregando -se com toda a sua emoção e subjetividade na atividade.
Durante a \textbf{brincadeira} a \textbf{criança} construirá o seu universo particular e terá a possibilidade de criar uma situação imaginária, na qual poderá assumir diferentes papéis e criar roteiros : como quando \textbf{brincam} de andar de ônibus e percebemos a divisão de papéis distribuídos e acordados pelas próprias \textbf{crianças} : o motorista, o trocador e os passageiros, cada qual com sua função ; e confrontar -se com dificuldades e angústias de forma criativa ; buscando soluções para resolvê -las ou mesmo liquidá -las em um espaço-tempo, onde podem experimentar o mundo e construir uma compreensão própria sobre as pessoas, os sentimentos e os diversos conhecimentos.
A situação imaginária é uma característica presente na \textbf{brincadeira} das \textbf{crianças} e vai sendo apreendida através da interação e \textbf{estímulos} ambientais.
Por exemplo, antes de aprender a \textbf{brincar}, uma espiga de milho é para a \textbf{criança} apenas um alimento e como tal é apreendido por ela, mas quando a imaginação passa a ocupar um papel importante na conduta \textbf{infantil}, uma espiga de milho pode vir a ser uma boneca e a \textbf{criança} se comporta em relação a ela de acordo com o novo sentido que ela lhe dá. ( VASCONCELOS, 2003, p. 10 ) O \textbf{brincar} também pode envolver a reiteração e a \textbf{criança} \textbf{demonstra} isso \textbf{repetindo} a situação até que se \textbf{satisfaça}, ou se sinta confiante.
Isso pode ocorrer também durante outras atividades prazerosas, como na leitura de uma história, por exemplo, quando solicitam que esta seja \textbf{contada} e recontada várias vezes.
Nesse sentido, Moyles ( 2002 ) afirma que o ato de \textbf{brincar} : > [... ] oferece situações em que as habilidades podem ser praticadas, tanto as físicas quanto as mentais, e \textbf{repetidas} tantas vezes quanto for necessário para a confiança e o domínio.
Além, disso, ele permite a oportunidade de explorar os próprios potenciais e limitações. ( p. 22 ) \textbf{Repetir} \textbf{brincadeiras} e jogos nem sempre estará ligado ao prazer, podendo estar também relacionada à necessidade que a \textbf{criança} tem em compreender algo novo ou angustiante que será \textbf{repetido} até que ela compreenda e \textbf{assimile} esta situação.
Ainda em relação ao aspecto da satisfação, a \textbf{criança} muitas vezes abre \textbf{mão} de seus \textbf{desejos} mais imediatos para continuar participando da \textbf{brincadeira}.
Basta observar no processo de negociação que acontece frequentemente entre as \textbf{crianças}, quando uma delas sugere algo que o grupo recusa e a \textbf{criança} tranquilamente diz : “ Tá bom! ” e continua a \textbf{brincar}.
O prazer, então, é alcançado na \textbf{manutenção} da atividade de \textbf{brincar} e não na realização imediata de \textbf{desejos}.
O jogo deve estar presente no conjunto das situações de aprendizagem e sempre que as \textbf{crianças} mostrem interesse em \textbf{brincar} com seus \textbf{pares}, o professor tem uma excelente oportunidade para observar e registrar como elas se organizam no grupo, como \textbf{brincam}, ou mesmo para observar uma \textbf{criança} que esteja lhe chamando a atenção.
São muitas as oportunidades que as \textbf{crianças} podem ter para \textbf{brincar} e aprender.
Algumas possibilidades de aprendizagem serão descritas a seguir : Os \textbf{bebês} \textbf{menores} de 1 ano podem aprender a : - \textbf{brincar} com os professores de cobrir e descobrir o rosto, de procurar e achar objetos \textbf{escondidos}, de \textbf{esconder} -se em algum canto da sala e ser encontrado ; - encaixar peças de madeira, empilhar cubos, puxar objetos etc. ; - envolver -se em troca de objetos com outros \textbf{bebês}, em ações ritmadas como \textbf{bater} as palmas das \textbf{mãos} sobre uma superfície, em entrar e sair de túneis etc. ; - \textbf{ouvir} e dançar, dentro de suas possibilidades, músicas e ritmos diversificados.
Com as \textbf{crianças} de 1 e 2 anos, novas aprendizagens podem ser estimuladas, tais como : - \textbf{brincar} de \textbf{roda} \textbf{imitando} \textbf{gestos} e cantos do professor e dos \textbf{colegas} ; - \textbf{brincar} de esconde-esconde, de jogar bola e de correr, com a supervisão do professor ; - \textbf{imitar} \textbf{gestos} e vocalizações de \textbf{adultos}, \textbf{crianças} ou animais, dentre outras, ou a \textbf{imitar} objetos ( o \textbf{som} do motor de um \textbf{carro}, por exemplo ) ; - \textbf{brincar} com diferentes ritmos e músicas variadas, vivenciando diversas possibilidades de \textbf{movimentação} corporal ( diferentes posições, velocidades, sensações ).
As \textbf{crianças} de 3 anos podem aprender a : - participar de \textbf{brincadeiras} de \textbf{roda}, sem precisar ter o professor como modelo ; - \textbf{brincar} de esconde-esconde e pega-pega, ou jogar bola, com supervisão do professor ; - \textbf{imitar} \textbf{gestos}, posturas e vocalizações de modelos ( \textbf{adultos}, \textbf{crianças}, animais ou \textbf{personagens} de histórias ) na ausência deles ; - utilizar objetos e \textbf{vestimentas} para assumir papéis no faz de conta onde reproduz situações cotidianas ; - reproduzir as ações de um \textbf{personagem} de uma história lida ( \textbf{imitar} o lobo da história, caminhar como os sete anões cantando na floresta ) ; - construir, ajudadas pelo professor, \textbf{brinquedos} com sucatas a partir de modelos, casas ou castelos com \textbf{areia}, sucata, tocos de madeira e outros materiais. \textbf{Crianças} de 4 anos podem aprender a : - comunicar -se com os companheiros na \textbf{brincadeira} utilizando \textbf{sons}, musicais ou não, ou diferentes formas de \textbf{gestos} e expressões vocais e corporais ; - \textbf{brincar} com a sonoridade de palavras, criando rimas ; - \textbf{brincar} de esconde-esconde, de jogar bola, de pique, de seguir o mestre, de lenço atrás, de caça ao tesouro etc. ; - \textbf{montar} quebra-cabeça ; - \textbf{dramatizar} uma história usando bonecos ou marionetes como atores ; - construir \textbf{brinquedos}, casas e cidades com diferentes materiais, sem necessariamente usar um modelo.
As \textbf{crianças} de 5 anos podem ser \textbf{apoiadas} a aprender, dentre outras coisas, a : - \textbf{brincar} de pula-sela, amarelinha, corda, pega-pega ; - elaborar coletivamente e \textbf{dramatizar} um enredo usando bonecos como atores ; - escolher \textbf{vestimenta} para compor um \textbf{personagem} para si ou para um \textbf{colega} ; - maquiar -se ou a um \textbf{colega} para desempenhar certo papel ; - discutir as intenções dos \textbf{personagens} de um enredo encenado ; - participar de jogos de tabuleiro como : loto, damas, memória, dominó etc. ; - explicar as regras de um jogo para outra \textbf{criança} ; - fazer \textbf{brinquedos} com sucata sem seguir modelo ; - construir casa/castelos de cartas, de cartolina, de panos e outros materiais ; - fazer dobraduras simples, elaborar máscaras, fazer bonecas de pano, ou de espiga de milho ; - construir e empinar pipas com a \textbf{ajuda} de um \textbf{adulto}. ( SEDUC, 2011, p. 41-42 ).
Essas são algumas das \textbf{brincadeiras} que podem ser organizadas diariamente.
Para tanto, a escolha do local e dos materiais que estarão disponíveis é fundamental.
Com relação à \textbf{brincadeira}, o professor precisa : - oferecer um repertório de cantigas, parlendas, adivinhas etc., possibilitando que as \textbf{crianças} vivenciem \textbf{brincadeiras} dançando, cantando, \textbf{imitando} ; - oportunizar situações em que as \textbf{crianças} possam \textbf{brincar} de faz de conta de diferentes formas : sozinhas, com o grupo, de forma livre e orientada pelo professor ; - estimular situações em que as \textbf{crianças} organizem enredos para as dramatizações, roteiros para a produção de danças e musicais e, ainda, planejem a confecção de \textbf{brinquedos} ; - respeitar o tempo e o ritmo das \textbf{crianças} enquanto \textbf{brincam} ; - mediar os conflitos surgidos nas \textbf{brincadeiras} ; - participar das \textbf{brincadeiras}, se solicitado.
A \textbf{brincadeira} é a atividade principal do \textbf{dia} a \textbf{dia} da \textbf{criança}.
Nas \textbf{brincadeiras}, as \textbf{crianças} podem tomar decisões, expressar sentimentos e valores, conhecer a si, aos outros e o mundo, \textbf{repetir} ações prazerosas, partilhar, expressar sua individualidade e identidade por meio de diferentes linguagens, usar o corpo, os sentidos, os movimentos, solucionar problemas e criar.
É possível fazer um rico exercício com a leitura da BNCC se buscarmos enxergar nos diversos Objetivos de Aprendizagem e Desenvolvimento situações diversas de interações das \textbf{crianças} e \textbf{brincadeira} que mobilizam suas aprendizagens e desenvolvimento.
Os Objetivos de Aprendizagem devem ser lidos sempre no contexto dos Campos de Experiências.
Tanto eles são possibilidades de aprendizagem em um Campo, como eles também atravessam outros Campos.
Por exemplo, se tomarmos o exemplo do jogo simbólico do Tomando Banho na Maré, veremos as \textbf{crianças} interagindo e construindo conhecimentos a partir da constituição de campos como : O eu, o outro e o nós ; Corpo, \textbf{gestos} e movimentos e \textbf{Escuta}, fala, pensamento e imaginação, e, a depender da reflexão, inclusive, outros campos.
Refletir, no contexto das interações reais acontecidas na turma de \textbf{crianças}, como os objetivos foram mobilizados no processo também pode enriquecer o olhar da(o) professora(r) para os processos educativos.
Vale aqui um destaque para o \textbf{cuidado} que devemos ter para não ler os objetivos de aprendizagem e desenvolvimento como objetivos pontuais a atingir ou não atingir.
Se olharmos atentamente para esses organizadores na BNCC, eles nos dizem de processos que vão se desenvolvendo continuadamente e de formas diversificadas.
Por exemplo, na \textbf{brincadeira}, quando as \textbf{crianças} usam \textbf{gestos} e movimentos para agregar novos elementos ao enredo, de certa forma elas criam com o corpo formas diversificadas de expressão, de sentimentos, sensações e emoções, elas mobilizam esses processos e seguem aprendendo sobre elas mesmas, sobre os \textbf{colegas} e suas formas de agir e pensar o mundo, assim como sobre o repertório de sua cultura.
São processos experienciados pelas \textbf{crianças} e pelos \textbf{adultos} que interagem com elas.

% matched lemmas: acolher, adulto, afeto, ajuda, apoiar, areia, assimilar, atento, bater, bebê, brincadeira, brincar, brinquedo, carro, colega, colo, conseguir, contar, conversar, crescimento, criança, cuidado, curiosidade, cópia, demonstrar, desejo, dia, dramatizar, esconder, escuta, escutar, estímulo, gesto, imitar, infantil, instituição, interessante, livro, manutenção, menino, montar, movimentação, mão, nascimento, ouvir, par, parceiro, partilha, pequeno, personagem, professora, pé, relacional, repetir, roda, satisfazer, sensorial, som, vestimenta, vontade
\end{textsample}
